\documentclass{article}
\usepackage[T1]{fontenc}
\usepackage{lmodern}
\usepackage{graphicx}
\usepackage{amsmath,amssymb}
\usepackage[a4paper,margin=2cm]{geometry}
\usepackage[absolute,overlay]{textpos}
\setlength{\TPHorizModule}{1mm}
\setlength{\TPVertModule}{1mm}
\usepackage[indonesian]{babel}
\usepackage{hyperref}
\setlength{\emergencystretch}{2em}
% unit: Halaman 68

\begin{document}

% === Halaman 68 (python/native) ===

2. Selanjutnya, ketiga register didetak selama 22 putaran tambahan. 

Selama 22 putaran tersebut, 22-bit dari nomor frame (Fn) dicampur 

dengan bit register berdasarkan skema berikut: 


Pada putaran 0  i < 22, 

                - bit Fn[i] ditambahkan ke bit LSB dari setiap register R dengan 

menggunakan operasi XOR:



R[0] = R[0]  Fn[i]

              - tiap register kemudian didetak


     Isi register pada akhir putaran menyatakan kondisi awal untuk    

     pembangkitan frame  sepanjang 228-bit. 

68

\end{document}
